% Generated draft: requires mathematical and source review.
\documentclass[11pt]{article}
\usepackage[margin=25mm]{geometry}
\usepackage{fontspec}
\setmainfont{Latin Modern Roman}
\usepackage{amsmath,amssymb,mathtools,unicode-math}
\setmathfont{Latin Modern Math}
\usepackage{enumitem,needspace,xurl,hyperref}
\setlength{\parindent}{0pt}
\setlength{\parskip}{5pt}
\setlength{\emergencystretch}{3em}
\newcommand{\N}{\mathbb{N}}
\newcommand{\Z}{\mathbb{Z}}
\newcommand{\Q}{\mathbb{Q}}
\newcommand{\R}{\mathbb{R}}
\newcommand{\C}{\mathbb{C}}
\newcommand{\F}{\mathbb{F}}
\newcommand{\A}{\mathbb{A}}
\newcommand{\PP}{\mathbb{P}}
\newcommand{\E}{\mathbb{E}}
\newcommand{\eps}{\varepsilon}
\newcommand{\dd}{\,\mathrm{d}}
\begin{document}
\section*{UnsolvedMath ID: 6200079}
\textbf{Source problem code:} catalogue entry
\subsection{Definitions and mathematical statement}
A finitely generated Coxeter group has a presentation
\[
 G=\langle S:(st)^{m_{st}}=1\text{ whenever }m_{st}<\infty\rangle,
\]
where $S$ is finite, $m_{ss}=1$, and for $s\ne t$ the symmetric entries $m_{st}$ lie in $\{2,3,\ldots\}\cup\{\infty\}$. Assume $G$ is Gromov hyperbolic, meaning its Cayley graph has uniformly thin geodesic triangles.

Let $\mathbb H^n$ denote real hyperbolic $n$-space and $\operatorname{Isom}(\mathbb H^n)$ its full isometry group with its usual topology. A discrete embedding is an injective homomorphism
\[
 \rho:G\longrightarrow\operatorname{Isom}(\mathbb H^n)
\]
whose image is a discrete subgroup. The source asks whether every Gromov-hyperbolic Coxeter group admits such an embedding for sufficiently large finite $n$; in particular, is
\[
 \forall G\text{ as above}\quad\exists n<\infty\quad
 \exists\rho\text{ a discrete embedding}
\]
true? The dimension may depend on $G$; no uniform dimension bound is stated. No torsion-free, orientation-preserving, cocompact, or convex-cocompact condition is added.

The source does not explicitly identify whether $\mathbb H^n$ means real hyperbolic space or another hyperbolic model, although real hyperbolic space is the usual reading of this notation. The research summary's parenthetical description of a discrete embedding as having a ``bounded orbit'' is not adopted: boundedness of an orbit is not the definition of discreteness and is a different dynamical property.
\subsection{Short English statement}
Can every hyperbolic Coxeter group act faithfully as a discrete group of isometries of some finite-dimensional real hyperbolic space? A linear representation alone does not supply the rank-one isometry target and discreteness requested here.

The supplied summary reports linear reflection representations and mentions Šapirovskii and Monod--Cornulier, but gives no exact theorem matching the target. It labels the question partial progress/open. The source's historical problem status and registry assessment are retained with current status explicitly needing review.
\subsection{Sources}
\begin{itemize}
\item Local UnsolvedMath registry, permanent ID 6200079, source code catalogue entry, ``Boundaries of Groups and Kleinian Groups --- Problem 79,'' proposed by Misha Kapovich, 2005. Dataset curation and metadata: Ulam AI and UnsolvedMath Contributors, CC BY 4.0.
\item Kapovich, \emph{Problems on Boundaries of Groups and Kleinian Groups} (2005), Problem 79, PDF page 20, \url{https://www.math.ucdavis.edu/~kapovich/EPR/problems.pdf}; normalized pdf-text extraction, accessed 2026-07-29.
\item The supplied summary names Šapirovskii and Monod--Cornulier for linearity-related work without complete references, exact statements, or a verified discrete-hyperbolic embedding theorem. Source-specific provenance marks current status and statement-reuse terms NEEDS\_REVIEW for the public author-hosted PDF, separately from the dataset licence.
\end{itemize}
\end{document}
